\begin{table}[!h]
\centering
\caption{\label{tab:s1_matching_double_counting_summary}Double counting between the AIS-based and non-broadcasting CO$_2$ estimates, 2017 to 2025, following the steps in the text. The share at risk is the total of Table \ref{tab:s1_matching_double_counting_by_gap}; because S1 images a vessel at an arbitrary moment, it is also the probability that a detection of a broadcasting vessel is missed by the matching. The matched detections are the ones that were not missed, which gives the number missed and their share of the unmatched detections. Non-broadcasting emissions are estimated from unmatched detections, so that share of them duplicates AIS-based emissions. The last two rows add the effect of the missed detections being absent from the matched detections the emissions regression model is trained on. The miss rate beyond 18 hours is taken to reach 100\% at 24 hours (Fig. \ref{fig:fig-s1-matching-double-counting}A, dashed).}
\centering
\fontsize{8}{10}\selectfont
\begin{tabular}[t]{>{\raggedright\arraybackslash}p{8cm}r}
\toprule
  & Value\\
\midrule
Share of AIS-based CO$_2$ at risk of double counting: emitted during AIS gaps in which a detection of the vessel would be missed & 5.1\%\\
Matched S1 detections & 17.0 million\\
Detections of broadcasting vessels missed by the matching & 0.92 million\\
Unmatched S1 detections & 14.9 million\\
Share of unmatched detections that are broadcasting vessels & 6.2\%\\
\addlinespace
Non-broadcasting CO$_2$ that duplicates AIS-based CO$_2$, 2017--2025 & 175 million tonnes (1.5\% of combined total)\\
Increase in the emissions per detection fitted by the emissions regression model, because the missed detections are absent from its training data & 5.4\%\\
Upper bound on non-broadcasting CO$_2$ that duplicates AIS-based CO$_2$, including that increase & 11.9\% of non-broadcasting CO$_2$ (2.9\% of combined total)\\
\bottomrule
\end{tabular}
\end{table}
